\chapter{Výsledky a vyhodnocení}
\label{ch:vysledky}

% Kapitola navazuje na implementaci (kap.~7) a vyhodnocuje navržený simulační model
% z~hlediska zadání: „Otestujte navrhnuté simulace. Vyhodnoťte výsledky simulací
% z~hlediska jejich dopadu na dopravu."
%
% Struktura kapitoly:
%   1. Ověření baseline modelu proti nezávislým datům (CSD, Waze).
%   2. Scénářové experimenty s reálnými uzavírkami z~databáze.
%   3. Křížová validace pomocí tabulek traffic_jams, accidents, event_links.
%   4. Citlivostní analýza parametrů modelu.
%   5. Diskuse zjištění a praktických implikací.
%   6. Omezení studie.


% ============================================================================
\section{Validace výchozího (baseline) stavu}
\label{sec:baseline}
% ============================================================================

% Cíl sekce: prokázat, že model bez uzavírek dostatečně věrně reprodukuje
% reálný dopravní stav -- je to nutný předpoklad pro důvěryhodnost scénářů.

\subsection{Shoda modelových intenzit s~celostátním sčítáním dopravy (CSD)}
% - Porovnat modelované objemy (link flows z~BFW přiřazení) s~pozorovanými
%   hodnotami AADT na kalibračních a zejména VALIDAČNÍCH profilech CSD 2020.
% - Kalibrační profily (65 \%) sloužily při ODME -- zde jen stručně rekapitulovat
%   dosažené metriky (R², \%RMSE, GEH, MAPE, bias, slope).
% - Hlavní důraz na NEZÁVISLOU validační sadu (35 \% profilů):
%   - scatter plot obs vs. mod s~regresní přímkou a 45° linií,
%   - distribuce GEH: kolik profilů splňuje GEH < 5 (standard: ≥ 85 \%),
%   - histogram reziduí (mod − obs) -- ověřit symetrii (žádný systematický bias).
% - Screenline analýza: součty toků přes definované řezy městem (mosty přes
%   Svratku, vnější/vnitřní okruh, radiály) -- porovnat s~pozorovanými součty.

\subsection{Validace rychlostí na síti pomocí Waze free-flow dat}
\label{subsec:freeflow}
% NOVÝ, ZAJÍMAVÝ PŘÍSTUP -- využití tabulky traffic_jams z~PostgreSQL.
%
% Myšlenka: Sloupec speed_normal_kmh v~záznamu o~zácpě udává rychlost
% za normálních podmínek na daném segmentu (volně plynoucí provoz).
% Agregací přes road_segments (→ osm_id → link_id) získáme nezávislý
% odhad volné rychlosti pro tisíce hran sítě.
%
% Postup:
%   1. Z~jams_segment_stats.parquet vzít median_speed_normal_kmh per segment.
%   2. Přes road_segments.osm_id namapovat na link_id v~AequilibraE síti.
%   3. Porovnat s~modelovou speed_ab / speed_ba (free-flow rychlost z~OSM + normalizace).
%
% Vizualizace:
%   - Scatter plot: Waze free-flow vs. model free-flow; R² a RMSE.
%   - Mapa odchylek: kde model přeceňuje / podceňuje rychlost.
%   - Rozlišit podle funkční třídy komunikace (motorway, primary, residential…).
%
% Význam: Na rozdíl od CSD (objemy) jde o validaci RYCHLOSTNÍ složky modelu --
% ta vstupuje do BPR funkce a přímo ovlivňuje cestovní časy a přesměrování.

\subsection{Prostorové vzorce kongesce: model vs.\ reálné zácpy}
\label{subsec:congestion-patterns}
% Porovnat modelové V/C poměry s~reálnou frekvencí zácp z~traffic_jams.
%
% Postup:
%   1. Pro každý road_segment spočítat počet zaznamenaných zácp (jam_count)
%      a průměrný delay_seconds -- z~jams_segment_stats.parquet.
%   2. Namapovat na link_id přes osm_id.
%   3. Porovnat s~modelovým V/C (volume / capacity z~baseline assignment).
%
% Očekávání: Hrany s~V/C > 0.8 by měly korelovat s~místy, kde jsou zácpy
% nejčastější. Pokud model predikuje nízké V/C tam, kde Waze reportuje
% časté zácpy, naznačuje to chybu v~kapacitě nebo poptávce.
%
% Vizualizace:
%   - Heatmapa překryvem: V/C z~modelu vs. hustota zácp z~Waze.
%   - Korelační graf segment-level: V/C vs. jam_count (Spearman ρ).
%   - Rozlišení špička / mimo-špičku pomocí časových agregací
%     z~traffic_jams (hodina dne, den v~týdnu).


% ============================================================================
\section{Experimenty se scénáři dopravních uzavírek}
\label{sec:scenare}
% ============================================================================

% Cíl: Otestovat model na reálných uzavírkách z~databáze restrictions a~vyhodnotit
% dopad na dopravní síť pomocí rozdílové analýzy (delta analysis) vůči baseline.

\subsection{Výběr testovacích uzavírek}
% - Z~restrictions vybrat 3--5 reprezentativních uzavírek s~různými charakteristikami:
%   a) Plná uzavírka páteřní komunikace (severity = full_closure, silnice I. třídy).
%   b) Částečná uzavírka s~redukcí pruhů (severity = lane_reduction).
%   c) Uzavírka v~centru města (vysoká hustota alternativ).
%   d) Uzavírka na periferní radiále (omezená redundance sítě).
%   e) Uzavírka mostu / přechodu přes vodní tok (topologicky kritický článek).
% - Pro každou uvést: link_id, road_number, city, restriction_type, severity,
%   valid_from/to, zda má liniovou geometrii a směrovost.

\subsection{Výsledky scénářového přiřazení}
% Pro každý vybraný scénář prezentovat:
%
% a) Síťové KPI:
%   - Celkový VHT (Vehicle Hours Traveled) a jeho změna ΔVHT vs. baseline.
%   - Celkový VKT (Vehicle Kilometers Traveled) a ΔVKT.
%   - Průměrný cestovní čas a jeho nárůst.
%   - Počet hran s~V/C > 1.0 (přetížení).
%
% b) Koridorová analýza:
%   - Mapy Δflow a Δtravel_time na dotčeném koridoru a jeho okolí.
%   - Identifikace objízdných tras, na které se doprava přesunula.
%   - Kvantifikace nárůstu zatížení na objízdných trasách (% nárůst).
%
% c) Kaskádové efekty:
%   - Kde se projevují sekundární kongesce (hrany vzdálené od uzavírky
%     s~nárůstem V/C > 0.2)?
%   - Jak daleko do sítě zasahuje dopad uzavírky (metricky: vzdálenost
%     nejbližší hrany se signifikantním Δflow od místa uzavírky)?

\subsection{Porovnání typů uzavírek}
% Napříč vybranými scénáři porovnat:
%   - Žebříček dle ΔVHT: která uzavírka má největší celosíťový dopad?
%   - Poměr lokálního vs. celosíťového dopadu.
%   - Nelineární efekty: při kombinaci dvou uzavírek současně (je-li
%     v~datech případ) -- je dopad superpozicí, nebo výrazně horší?
% - Tabulka: scénář | typ | ΔVHT | ΔTT | max ΔV/C | počet přetížených hran.

\subsection{Vliv liniové geometrie a~směrovosti na výsledky}
\label{subsec:line-vs-point}
% NOVÝ EXPERIMENT -- porovnání dvou režimů matchingu uzavírek:
%
% a) Bodový matching (starý přístup):
%   - Uzavírka přiřazena k~nejbližší hraně dle bodu (lat, lon).
%   - Riziko: u~dlouhých uzavírek zasáhne jen jednu hranu.
%
% b) Liniový matching (nový přístup):
%   - Uzavírka s~line_wkt bufferována a průnik s~hranami sítě.
%   - Zasáhne všechny hrany, které geometricky protíná.
%
% c) Směrový vs. obousměrný režim:
%   - closure_direction = aligned/opposite → aplikace jen na ab/ba.
%   - vs. výchozí both → obousměrná redukce.
%
% Metriky pro porovnání:
%   - Počet zasažených link_id (bodový vs. liniový).
%   - Rozdíl v~ΔVHT a ΔTT pro tentýž scénář.
%   - Případové studie: uzavírka dlouhé rekonstrukce na silnici I. třídy,
%     kde bodový matching zasáhne 1 hranu, ale liniový 15+ hran.
%   - Vizualizace: mapa zasažených hran při obou přístupech.


% ============================================================================
\section{Křížová validace modelu vůči nezávislým datovým zdrojům}
\label{sec:krizova-validace}
% ============================================================================

% Cíl: Ověřit predikce modelu pomocí dat, která nebyla použita při kalibraci.
% Využít bohaté datové zdroje z~PostgreSQL (traffic_jams, accidents,
% event_links, alerts) jako nezávislé referenční body.

\subsection{Validace dopadů uzavírek pomocí kauzálních vazeb (event\_links)}
\label{subsec:event-links-validace}
% KLÍČOVÝ A ORIGINÁLNÍ VALIDAČNÍ PŘÍSTUP.
%
% Tabulka event_links mapuje: restriction_id → traffic_jam_id + delay_seconds.
% Pro vybrané uzavírky tedy ZNÁME:
%   - Kolik zácp reálně způsobily (jams_caused).
%   - Na kterých segmentech se zácpy objevily (→ link_id přes road_segments).
%   - Celkové reálné zpoždění (total_delay_s z~restriction_impact.parquet).
%
% Validační postup:
%   1. Vybrat uzavírky z~experimentů (sec. 8.2), pro které existují event_links.
%   2. Spustit scénář v~modelu → získat ΔV/C a Δtravel_time na všech hranách.
%   3. Porovnat:
%      a) PROSTOROVĚ: Na kterých hranách model predikuje kongesci (ΔV/C > práh)?
%         Shodují se s~hranami, kde Waze zaznamenal zácpy?
%         → Precision / Recall analýza: kolik reálných zácp model „odchytil".
%      b) KVANTITATIVNĚ: Celkový modelový nárůst VHT vs. reálné total_delay_s.
%         → Scatter plot přes více uzavírek; regrese a R².
%      c) TOPOLOGICKY: Leží modelové kongesce na stejných alternativních
%         trasách jako reálné zácpy, nebo model predikuje přesměrování jinam?
%
% Očekávání a interpretace:
%   - Statický model zachytí hlavní přelivy (primární objízdné trasy), ale
%     nemusí zachytit kaskádové zácpy (ty vyžadují dynamický model).
%   - I~částečná shoda (>50 \% recall) je pro statický model dobrý výsledek
%     a potvrzuje praktickou použitelnost.

\subsection{Korelace přetížení modelu s~nehodovými blackspoty}
\label{subsec:accident-correlation}
% Myšlenka: Úseky s~vysokým poměrem V/C v~modelu indikují vyšší míru
% kongesce a tedy potenciálně vyšší nehodovost (rear-end kolize, nehody
% při řazení, stresové chování řidičů).
%
% Postup:
%   1. Z~accidents agregovat počet nehod per road_segment (→ osm_id → link_id).
%   2. Modelové V/C z~baseline přiřazení per link.
%   3. Spearman rank korelace: V/C vs. accident_count per link.
%   4. Logistická regrese: P(nehoda na linku) ~ V/C + road_type + speed_limit.
%
% Rozšíření:
%   - Rozlišit nehody za mokra (weather_condition = rain/wet) vs. suché podmínky.
%   - Porovnat s~alerts typu HAZARD -- zda alerty kopírují modelové přetížení.
%
% Vizualizace:
%   - Mapa: překryv heatmapy nehod s~V/C > 0.9 z~modelu.
%   - Box plot: distribuce accident_count pro kvartily V/C.
%
% Význam: Nejde o primární účel dopravního modelu, ale tato korelace dodává
% důvěryhodnost -- pokud model správně identifikuje přetížené úseky, měl by
% se to promítnout i do nezávislých dat o nehodovosti.

\subsection{Porovnání modelových a~reálných kongescí v~časových profilech}
\label{subsec:temporal-profiles}
% Statický model AequilibraE počítá denní průměr. Porovnání s~časovými profily
% z~traffic_jams umožňuje posoudit, zda model odpovídá „špičkovému" chování.
%
% Postup:
%   1. Z~traffic_jams agregovat: pro každý segment průměrný počet zácp
%      a~delay per hodina dne (6--9, 9--15, 15--18, 18--22).
%   2. Model produkuje denní V/C. Porovnat, zda modelový V/C koreluje
%      se špičkovou (7--9, 15--17) frekvencí zácp.
%
% Hypotéza: Pokud model dobře zachycuje prostorové rozložení poptávky,
% pak by V/C > 0.8 mělo korelovat se segmenty, kde jsou zácpy nejčastější
% ve špičkových hodinách.
%
% Doplňkový experiment: Pokud bude implementován learn-profile (krok 18 pipeline),
% porovnat výstup časového profilu modelu s~hodinovou distribucí zácp z~Waze.


% ============================================================================
\section{Citlivostní analýza}
\label{sec:citlivost}
% ============================================================================

% Cíl: Ověřit robustnost zjištění při změnách klíčových vstupních parametrů.

\subsection{Citlivost na úroveň poptávky}
% - Spustit baseline a vybrané scénáře s~O-D maticí škálovanou na 90\,\%
%   a 110\,\% původní poptávky.
% - Sledovat: Jak se mění ΔVHT, ΔTT a prostorový vzorec přesměrování?
% - Očekávání: Při vyšší poptávce budou dopady uzavírek nelineárně horší
%   (síť je blíže kapacitnímu limitu → menší rezerva pro přesun dopravy).
% - Vizualizace: graf ΔVHT vs. úroveň poptávky pro vybrané scénáře.

\subsection{Citlivost na parametry severity uzavírky}
% - Pro stejnou uzavírku variovat capacity_factor (0.0, 0.1, 0.3, 0.5)
%   a~speed_factor (0.3, 0.5, 0.7, 1.0).
% - Sledovat nelineární práh: od jakého snížení kapacity dopad skokově roste?
% - Porovnat plnou uzavírku (full) vs. redukci na 1 pruh (lanes) vs.
%   pouze snížení rychlosti (speed penalty).

\subsection{Citlivost na quality\_score filtraci uzavírek}
% - Spustit scénář se všemi uzavírkami (min_quality_score = 0) vs.
%   jen s~kvalitními (min_quality_score ≥ 70).
% - Zjistit: Kolik nekvalitních uzavírek zbytečně zatěžuje model?
%   Mění se výrazně KPI při filtraci?

\subsection{Vliv konvergenčních parametrů přiřazení}
% - Variovat RGAP toleranci (1e-3, 1e-4, 1e-5) a maximální iterace BFW.
% - Sledovat stabilitu výsledných Δ-metrik: pokud se ΔVHT liší o~méně
%   než 1\,\% při zpřísnění konvergence, model je dostatečně konvergovaný.


% ============================================================================
\section{Diskuse zjištění}
\label{sec:diskuse}
% ============================================================================

\subsection{Kritické články sítě a~odolnost}
% - Na základě experimentů identifikovat hrany, jejichž uzavírka způsobuje
%   největší ΔVHT -- tyto představují kritické články sítě (bottlenecks).
% - Diskutovat redundanci: oblasti s~více alternativami (centrum, mřížková
%   zástavba) vs. oblasti s~jedinou přístupovou cestou (mosty, periferní radiály).
% - Souvislost s~konceptem network resilience z~literatury (Sheffi, Cascetta).

\subsection{Přesnost statického modelu pro hodnocení uzavírek}
% - Na základě křížové validace (sec.~\ref{sec:krizova-validace}) diskutovat:
%   - Kde statický model funguje dobře (hlavní přelivy, denní průměry)?
%   - Kde selhává (dynamické kaskády, špičková kongesce, krátkodobé incidenty)?
% - Porovnat s~teoretickými limity statického přiřazení z~kapitoly 2.
% - Navrhnout, kdy je nutné přejít na DTA (dynamické přiřazení).

\subsection{Přínos rozšířených dat pro kvalitu simulace}
% - Kvantifikovat přínos liniové geometrie vs. bodového matchingu
%   (z~experimentu \ref{subsec:line-vs-point}): o~kolik se zlepšily KPI?
% - Přínos směrovosti: Diskutovat případy, kde ignorování směru vedlo
%   k~nadhodnocení dopadu (uzavírka jednoho směru zbytečně penalizuje oba).
% - Přínos Waze free-flow rychlostí (sec.~\ref{subsec:freeflow}) pro kalibraci:
%   umožňují zpřesnit speed_ab/ba na tisících hran, kde OSM data chybí.

\subsection{Praktické implikace pro dopravní plánování}
% - Jak mohou výsledky sloužit dopravním inženýrům při plánování uzavírek:
%   - Identifikace nejhorších scénářů ještě před zahájením stavby.
%   - Odhad objízdných tras a jejich potřebné kapacity.
%   - Informační strategie pro řidiče (které alternativy doporučit).
% - Možnost integrace s~řízením SSZ (světelná signalizace) na objízdných trasách.
% - Vazba na práci Ondruškové: jak vizualizace dopadů v~mapě pomáhá
%   při komunikaci s~veřejností a stakeholdery.


% ============================================================================
\section{Omezení studie}
\label{sec:omezeni}
% ============================================================================

\subsection{Omezení modelu}
% - Statický přístup: nepodchycuje dynamiku v~čase (ranní náběh, rozptyl
%   po otevření). Vhodný pro strategické hodnocení, ne pro operativní řízení.
% - Chybějící módní volba: model pracuje pouze s~individuální automobilovou
%   dopravou. Přesun cestujících na MHD/kolo není modelován.
% - Zjednodušení křižovatek: AequilibraE nemodluje detaily křižovatek
%   (signální plány, kapacity odbočení) -- lokální kongesce v~uzlech
%   může být podhodnocena.
% - Homogenní řidiči: v~rámci každé třídy (local/through) identické chování;
%   heterogenní VOT a informovanost nejsou zohledněny.

\subsection{Omezení dat}
% - Waze data (traffic_jams, alerts): pokrytí závisí na penetraci aplikace.
%   Venkovské oblasti a méně frekventované komunikace mohou být podreprezentovány.
% - Sloupec quality_score v~restrictions: 12\,\% záznamů má NULL.
%   Data NDIC/DATEX~II mohou obsahovat duplicity nebo zastaralé záznamy.
% - Accidents: pouze PČR hlášené nehody, nezachycují drobné kolize bez policie.
% - CSD profily: cca 400 profilů pro Brno -- pokrytí sítě není rovnoměrné,
%   menší komunikace mohou být pod-reprezentovány.
% - O-D matice: odvozena ze SLDB 2021 dojížďkových dat -- odráží denní
%   dojížďku, ne veškerou mobilitu (nákupy, volný čas, nákladní doprava
%   je řešena pouze odhadem).

\subsection{Přenositelnost a zobecnitelnost}
% - Metodika (pipeline, scénářové porovnání, validační přístupy) je obecná
%   a přenositelná na jiná města (viz init-city mechanismus, kap.~7).
% - Konkrétní kalibrační výsledky a zjištění o~kritických článcích platí
%   specificky pro Brno a nelze je přímo zobecnit.
% - Doporučení pro aplikaci na jiné město: minimální požadavky na data
%   (OSM síť, sčítací profily, O-D seed), ideálně doplnit o~Waze data.
